\begin{abstract}
임피던스 제어는 로봇이 외부 환경과 접촉할 때 위치를 강제로 추종하는 대신, 힘과 운동 사이의 동역학적 관계를 설계하는 제어 방법이다.
Physical AI와 로봇 기초모델 논의가 실제 조작 문제로 옮겨가려 할수록, 입문자는 마지막 접촉 단계에서 힘과 운동의 관계를 이해해야 한다.
물체, 사람, 작업물과 안정적으로 맞닿아 일하려면 반복 위치 추종이나 위치 정밀도만으로는 부족하며, 접촉력, 순응성, 강성, 감쇠를 함께 다루는 힘/임피던스 관점이 필요하다.
하지만 임피던스라는 용어는 전기회로, 기계진동, 로봇제어에서 서로 다른 표기와 직관으로 등장하기 때문에 처음 학습하는 입장에서는 개념의 연결이 잘 보이지 않는다.
이 논문은 비전공자와 로봇 제어 입문자가 임피던스를 입력-출력 및 에너지 관점에서 읽을 수 있도록 쓴 리뷰/튜토리얼이다.
저항만으로 부족했던 역사적 문제에서 출발해, 같은 언어로 비교하기 위한 LTI 시스템, 저장과 소산을 보여주는 2차 전기/기계 시스템, 다관절 로봇의 관절공간과 작업공간 문법, 로봇 임피던스 제어로 이어지는 학습 경로를 정리한다.
MuJoCo 실험실은 정밀 디지털 트윈이 아니라, 질량-스프링-댐퍼, 2자유도 평면 매니퓰레이터, Franka Panda 모델의 위치 액추에이터 기반 가상 벽/목표점 후퇴 실험으로 이 관계를 확인하는 학습 도구이다.
각 실험은 YAML 설정, CSV/NPZ 로그, summary JSON, 필요할 때 저장한 플롯 이미지, worksheet 메모를 반복 가능한 학습 기록으로 남기며, 학습자는 예측하고 관찰하고 다음 실험을 고르는 과정을 반복한다.
여기서 비교하려는 것은 특정 제어기의 성능 순위가 아니라, 전기-기계-로봇 시스템에서 반복해서 나타나는 힘과 운동의 해석 방식이다.
독자는 최종적으로 강성, 감쇠, 관성의 변화가 접촉력, 위치 오차, 오버슈트, 정착시간에 어떤 영향을 주는지 말할 수 있어야 한다.
\end{abstract}
